\subsection*{3.2. 偶相对维数}

保留 3.1 的记号, 并取偶数 $n=2m$. 为简便起见, 假设 $S$ 严格 Hensel(一般情形可由下降得到).

\subsubsection*{命题 3.2.1.}
\begin{enumerate}
\item[(i)] \underline{对 $x\in\Sigma$, 群 $H^n_{\{x\}}(X_s,R\psi_{\bar\eta}(\Lambda(m)))$ 有一个在正负号意义下确定的自然生成元 $\delta$, 它由关于 $\Lambda$ 的自然性及以下条件刻画:}
\[
(\delta,\delta)=(-1)^m\,2.
\]
\underline{当 $n=0$ 时, 再加上条件 $\operatorname{Tr}(\delta)=0$.}

\item[(ii)] \underline{对惯性群 $I$ 的某个与 $\Lambda$ 无关的特征标}
\[
\varepsilon_x:I\longrightarrow\{\pm1\}
\]
\underline{变差为}
\[
\operatorname{Var}(\sigma)(a)
=(-1)^m\left(\frac{\varepsilon_x(\sigma)-1}{2}\right)
(a,\delta)\delta
\]
\underline{其中 $a\in R^n\Phi_{\bar\eta}(\Lambda(m))$. 因而,}
\[
\sigma(\delta)=\varepsilon_x(\sigma)\delta.
\]
\end{enumerate}
\footnote{印本在 (i) 中写为 $H^n_{\{x\}}(R^n\psi_{\bar\eta}(\Lambda(m)))$. 证明明确使用了 2.2.5(D), 其中自然生成元属于 $H^n_{\{x\}}(X_s,R\psi_{\bar\eta}(\Lambda(m)))$.}

取逆极限后, 只需处理普遍情形 $\Lambda=\mathbb Z_\ell$. 此时在相差正负号的意义下, 只有一个 $\delta$ 能满足 (i) 的条件. 如 3.1.2 中那样归结为情形 2.2.5(*), 再应用 2.2.5(D).

\subsubsection*{补充 3.2.2.}
假设在 $x\in\Sigma$ 处, $X$ 在 $x$ 处的 Hensel 化 $X_{(x)}$ 同构于由方程
\[
\sum_{i\leq j}a_{ij}x_ix_j=0.
\]
定义的射影概形 $Q$ 在 $x_0$ 处的 Hensel 化. 则特征标 $\varepsilon_x$ 由分式域的可分二次扩张 $Z(C^+(Q))$ 定义. 这由 2.2.5(D) 得出.

\subsubsection*{3.2.3.}
当剩余特征不等于 $2$ 时, 补充 3.2.2 可更加显式. 此时 $I$ 有唯一一个阶为 $2$ 的非平凡特征标 $\varepsilon$. 对一个一致化元 $t$, 有
\[
\sigma(\sqrt t)=\varepsilon(\sigma)\sqrt t.
\]
另一方面, 此时 $X_{(x)}$ 同构于仿射概形
\[
\sum a_{ij}x_ix_j=b,
\]
在 $0$ 处的 Hensel 化, 其中 $b$ 属于极大理想, $Q(x)=\sum a_{ij}x_ix_j$ 非退化. 理想 $(b)$ 定义 $X_{(x)}$ 中作为 $J^n_{X_{(x)}/S}$ 支集的子概形的像. Clifford 代数的中心同构于
\[
k(\eta)\left(\left((-1)^{m+1}2b\det(a_{ij})\right)^{1/2}\right)
\]
(Bourbaki, \emph{Alg\`ebre}, 第 9 章, \S9, no.~4, 最后的注), 故
\[
\varepsilon_x=\varepsilon^{v(b)}.
\]
若 $b$ 的赋值为偶数, 则变差为零; 否则 $\varepsilon_x=\varepsilon$.